% Generated draft: requires mathematical and source review.
\documentclass[11pt]{article}
\usepackage[margin=25mm]{geometry}
\usepackage{fontspec}
\setmainfont{Latin Modern Roman}
\usepackage{amsmath,amssymb,mathtools,unicode-math}
\setmathfont{Latin Modern Math}
\usepackage{enumitem,needspace,xurl,hyperref}
\setlength{\parindent}{0pt}
\setlength{\parskip}{5pt}
\setlength{\emergencystretch}{3em}
\newcommand{\N}{\mathbb{N}}
\newcommand{\Z}{\mathbb{Z}}
\newcommand{\Q}{\mathbb{Q}}
\newcommand{\R}{\mathbb{R}}
\newcommand{\C}{\mathbb{C}}
\newcommand{\F}{\mathbb{F}}
\newcommand{\A}{\mathbb{A}}
\newcommand{\PP}{\mathbb{P}}
\newcommand{\E}{\mathbb{E}}
\newcommand{\eps}{\varepsilon}
\newcommand{\dd}{\,\mathrm{d}}
\begin{document}
\section*{UnsolvedMath ID: 6700036}
\textbf{Source problem code:} catalogue entry
\subsection{Definitions and mathematical statement}
Let \(\Gamma\) be a discrete countable group, let \(B=B\Gamma\) be a classifying space, and let \(f:X\to B\) be continuous with \(X\) a Riemannian manifold. The source first asks for a compact subset \(B_0\subset B\) containing \(f(X)\), and then asks for arbitrarily small curvature in nonzero integer multiples of an indicated pullback K-theory class.

If \(L_0\) is the intended fixed bundle and all categories are specified, its displayed target has the quantifier pattern
\[
 \exists B_0\subset B\text{ compact},\ f(X)\subseteq B_0:\quad
 \forall\varepsilon>0\ \exists N\in\mathbb Z\setminus\{0\}\ 
 \exists(L,\nabla)\text{ over }X:
 \quad\|\operatorname{curv}(\nabla)\|\le\varepsilon,
 \quad [L]=Nf^*[L_0].
\]
The integer and bundle may depend on \(\varepsilon\); no single fixed multiple for all errors is required by the excerpt. For positive \(N\), the usual finite-dimensional Whitney sum has class \([f^*L_0^{\oplus N}]=Nf^*[L_0]\). Negative \(N\) is meaningful as an additive K-theory class, not as an ordinary negative Whitney sum, unless an appropriate virtual/Fredholm model is supplied.

The source does not define Fredholm bundle, Fredholm coarea, K-theoretic equivalence, or the curvature norm. A common finite-rank curvature norm would use the supremum of the fiber-operator norm on unit tangent two-planes, but the intended Fredholm setting may differ. The bundle \(L_0\), its base \(B\) or \(B_0\), and whether its class is given, arbitrary, or existential are also missing. In particular a nonzero integer does not by itself require its multiple of a K-class to remain nonzero; that additional nontriviality must not be silently imposed.

No compactness hypothesis on \(X\) or relative compactness of \(f(X)\) is stated, so existence of the initial compact \(B_0\) is not an automatic consequence of the supplied assumptions. The classifying-space model and any required finite subcomplex condition need confirmation. These gaps prevent the displayed conditional pattern from being a fully specified theorem.

The summary describes this as an open Fredholm K-area/coarea question related to almost-flat bundles, Novikov theory, and scalar-curvature obstructions, with no dedicated resolution located. No precise auxiliary theorem or definition is supplied, and no independently verified status claim is made.
\subsection{Short English statement}
Can suitable nonzero multiples of the intended pullback K-class be represented by Fredholm bundles of arbitrarily small curvature, after locating the image in a compact part of \(B\Gamma\)? The compactness assumption, bundle quantifier, Fredholm definitions, and treatment of negative multiples are missing.
\subsection{Sources}
\begin{itemize}
\item Local registry: \url{data/unsolvedmath/problems.json}, permanent UnsolvedMath ID \(6700036\), source code catalogue entry; supplied statement, background, and research summary, updated 2026-08-06.
\item Misha Gromov, \emph{101 Questions, Problems and Conjectures around Scalar Curvature} (2017), printed tag [?39], occurrence 38, PDF page 45: \url{https://www.ihes.fr/~gromov/wp-content/uploads/2018/08/101-problemsOct1-2017.pdf}. Supplied normalized PDF-text extraction accessed 2026-07-29; described as incomplete.
\item The summary mentions almost-flat bundles and Novikov-related work without supplying exact definitions, original papers, or a resolution reference.
\end{itemize}
\end{document}
